% =====================================================================
%  Основное содержимое отчёта.
% =====================================================================

\labpart{Краткое описание разработанного приложения}

В backend (\texttt{SecureApp.Api}) добавлен класс
\texttt{InMemoryRecordStore} -- синглтон-сервис, хранящий записи в двух
\texttt{ConcurrentDictionary} (для конфиденциальных и неконфиденциальных
данных) без единого обращения к Entity Framework Core, файловой системе
или иному постоянному хранилищу; идентификаторы записей выдаются
атомарно через \texttt{Interlocked.Increment}, что делает сервис
безопасным для параллельных запросов без явных блокировок. Поверх сервиса
добавлены два REST-контроллера,
\texttt{InMemoryConfidentialController} и \texttt{InMemoryPublicController},
идентичные по контракту (маршруты, авторизация, разбиение по владельцу)
контроллерам постоянных данных из лабораторной работы №2.1, но
обслуживающие маршруты \texttt{/api/memory/confidential} и
\texttt{/api/memory/public}. На фронтенде добавлены две вкладки
«In-memory confidential» и «In-memory public», использующие тот же
компонент таблицы с поиском, что и постоянные данные, -- различие видно
пользователю только по надписи «RAM only -- not persisted, lost on
restart» на соответствующей вкладке.

Конфиденциальные данные шифруются алгоритмом AES-256-GCM (тем же классом
\texttt{ConfidentialCrypto}, что и в лабораторной работе №2.1)
непосредственно в момент создания или обновления записи -- то есть в
структурах данных самого сервиса никогда не хранится ни одной
конфиденциальной строки в открытом виде, только зашифрованные байты
(\autoref{lst:secureapp-memstore}).

\begin{codelisting}[fontsize=\footnotesize]{csharp}{Создание записи: шифрование выполняется до попадания данных в словарь}{secureapp-memstore}
public ConfidentialEntry CreateConfidential(int ownerId, string title, string content)
{
    var id = Interlocked.Increment(ref _nextConfidentialId);
    var now = DateTime.UtcNow;
    var entry = new ConfidentialEntry(id, ownerId, Encrypt(title), Encrypt(content), now, now);
    _confidential[id] = entry;
    return entry;
}
\end{codelisting}

\labpart{Анализ защищённости приложения с точки зрения кода}

\subsection*{1. Обработка невалидных данных}

Семь видов заведомо некорректного ввода были поданы на новый маршрут
напрямую (см. \autoref{lst:secureapp-invalid-input}): значения не того
типа (числа и булево значение вместо строк), синтаксически некорректный
JSON, отсутствие обязательного поля, превышение ограничения длины поля
(10\,000 символов), управляющие символы (перевод строки, символ
табуляции) внутри значения и полностью пустое тело запроса. Все семь
случаев отклонены кодом \texttt{400 Bad Request} системой валидации
ASP.NET Core, кроме управляющих символов внутри значения -- они являются
полноправным содержимым текстового поля, а не структурной ошибкой, и
корректно приняты (\texttt{201 Created}). После всех проверок
контрольный запрос к тому же маршруту подтвердил, что процесс не
аварийно завершился и не завис -- ни одно из некорректных значений не
привело к необработанному исключению на уровне процесса, а осталось
изолированным в рамках обработки одного HTTP-запроса.

\codelistingfile{text}{lab/code/invalid_input_output.txt}{Реакция на некорректный ввод: 7 проверок, 0 падений процесса}{secureapp-invalid-input}

\subsection*{2. Работа с конфиденциальными данными в коде}

В методы \texttt{ConfidentialCrypto.Encrypt}/\texttt{Decrypt} добавлена
явная очистка собственных промежуточных буферов
(\texttt{Array.Clear(plainBytes)}) сразу после использования, в блоке
\texttt{finally} -- то есть даже при исключении внутри самого
шифрования байтовый буфер не останется висеть в памяти дольше
необходимого. Это осознанно называется в коде и в этом отчёте \emph{мерой
«наилучшего возможного» (best-effort)}, а не гарантией: строка
\texttt{string} в .NET неизменяема, и исходная строка с открытым текстом,
переданная методу вызывающим кодом (в том числе десериализатором JSON,
создавшим её из тела HTTP-запроса), никаким программным способом самим
методом очищена быть не может -- он контролирует только свои
собственные, дополнительно созданные буферы. Управляемая среда
исполнения (.NET CLR) в этом отношении принципиально отличается от кода
на C/C++, где буфер с паролем можно явно обнулить сразу после
использования и быть уверенным, что других копий не осталось: в
управляемой среде таких гарантий нет, и, как показано далее в разделе
«Анализ дампов оперативной памяти», это не абстрактная угроза, а
наблюдаемое на практике поведение. По этой же причине класс
\texttt{System.Security.SecureString} сознательно не используется:
начиная с .NET Core он не даёт защиты вне Windows и, что важнее, не
защищает от ровно той угрозы, которая здесь актуальна -- копий строки,
уже созданных инфраструктурой ASP.NET Core до попадания данных в код
приложения.

\subsection*{3. Потокобезопасность}

Поскольку все операции доступны нескольким одновременным HTTP-запросам,
хранилище построено на \texttt{ConcurrentDictionary} вместо обычного
\texttt{Dictionary} с ручной блокировкой, а выдача идентификаторов -- на
атомарном \texttt{Interlocked.Increment} вместо чтения и увеличения
обычного поля, что исключает состояние гонки (два запроса, получившие
один и тот же идентификатор) без единого явного \texttt{lock} в коде
сервиса.

\labpart{Анализ оперативной памяти и процессорного времени приложения}

Оперативная память и загрузка процессора процесса backend измерялись
инструментом \texttt{dotnet-counters} (аналог \texttt{jvisualvm} для
.NET) в двух режимах: в состоянии покоя и под нагрузкой в 900
последовательных HTTP-запросов (300 циклов
создание~$\to$~обновление~$\to$~удаление записи) за приблизительно 10
секунд.

\begin{itemize}
  \item \textbf{Резидентная память (working set).} В состоянии покоя --
  около 259~МБ; на пике нагрузки -- около 293~МБ (прирост
  $\approx$34~МБ на 900 операций), что подтверждает: временные объекты
  создаются и достаточно быстро собираются сборщиком мусора, без
  неограниченного роста памяти.
  \item \textbf{Процессорное время.} В состоянии покоя процесс потреблял
  около 0,13~с пользовательского времени CPU на секунду
  наблюдения ($\approx$13\,\% одного ядра); под нагрузкой -- до
  1,09~с на секунду (то есть свыше одного полностью загруженного
  ядра, за счёт нескольких потоков одновременно), возвращаясь к уровню
  покоя сразу после завершения нагрузки.
  \item \textbf{Скорость выделения памяти в куче (allocation rate).} От
  0,7 до 1,7~МБ в секунду под нагрузкой -- ожидаемо для рабочей
  нагрузки, где на каждый HTTP-запрос выделяются короткоживущие объекты
  (строки из JSON, объекты записей).
  \item \textbf{Пропускная способность пула потоков.} До 877
  завершённых рабочих элементов в секунду -- согласуется с 900
  запросами, реально обработанными за $\approx$10 секунд.
  \item \textbf{Паузы сборщика мусора.} На гранулярности в 1 секунду не
  зафиксировано значимых пауз сборщика мусора -- нагрузка данного
  масштаба не вызывает заметных для пользователя задержек ответа
  из-за GC.
\end{itemize}

Совокупно приложение демонстрирует управляемый рост потребления ресурсов
под нагрузкой без признаков утечки памяти или чрезмерной загрузки
процессора для нагрузки такого масштаба.

\labpart{Анализ дампов оперативной памяти}

Дампы снимались утилитой \texttt{dotnet-dump collect} (аналог
\texttt{jstack}/полного дампа кучи JVM) в трёх контрольных точках,
предписанных заданием: после создания, после обновления и после удаления
одной конфиденциальной и одной неконфиденциальной записи, каждая с
уникальным маркером в содержимом (например,
\texttt{CONFIDENTIAL-PLAINTEXT-MARKER-AXQ7391}). Перед третьим дампом
дополнительно принудительно вызывалась полная сборка мусора -- побочный
эффект команды \texttt{dotnet-gcdump collect}, которая перед снятием
дампа кучи запрашивает у среды выполнения блокирующую сборку, -- чтобы
проверить, действительно ли удаление записи из структуры данных
приложения приводит к исчезновению значения из физической памяти
процесса после того, как сборщик мусора получил шанс отработать. Каждый
дамп затем анализировался утилитой \texttt{strings} с поиском обоих
маркеров -- методика, аналогичная поиску секретов в файле базы данных
из лабораторной работы №2.1, но применённая к дампу процесса, а не к
файлу на диске.

\codelistingfile[fontsize=\footnotesize]{text}{lab/code/memory_dump_findings.txt}{Число вхождений маркеров в каждом из трёх дампов}{secureapp-memory-dumps}

Результат оказался одинаковым во всех трёх контрольных точках и для
обеих версий значения (до и после обновления): и конфиденциальный, и
неконфиденциальный маркер найдены в дампе процесса по три раза каждый --
включая дамп, снятый после удаления записи и принудительной полной
сборки мусора. Иначе говоря, ни шифрование данных перед сохранением в
структуру приложения, ни последующее удаление записи, ни явно
запрошенная сборка мусора не убрали открытый текст конфиденциального
значения из памяти процесса.

Это не свидетельствует об ошибке именно в коде хранилища
(\texttt{InMemoryRecordStore} нигде не хранит открытый текст -- в его
собственных структурах данных лежат только зашифрованные байты, что было
подтверждено отдельно на уровне кода в предыдущем разделе). Наиболее
вероятное объяснение, согласующееся с внутренним устройством ASP.NET
Core: конвейер обработки HTTP-запроса (Kestrel и сериализатор
\texttt{System.Text.Json}) для производительности активно использует
переиспользуемые буферы (\texttt{ArrayPool<byte>}/\texttt{ArrayPool<char>}
из общего пула), в которые копируется тело запроса при разборе JSON и
тело ответа при его формировании. Буфер, возвращённый в пул, не
обнуляется и с точки зрения сборщика мусора не является мусором вовсе --
на него по-прежнему ссылается сам пул, поэтому произвольно долго
переживает не только удаление записи из словаря приложения, но и
принудительную полную сборку мусора. Тот факт, что оба маркера --
предназначенный для шифрования и предназначенный для хранения в открытом
виде -- обнаружены с одинаковой кратностью (по три вхождения),
подтверждает эту гипотезу: искать нужно не в структурах данных
приложения (там разница есть -- она в шифровании), а в разделяемой
инфраструктуре, которая одинаково обрабатывает оба вида данных на пути
от сети до кода приложения и обратно.
